% Options for packages loaded elsewhere
\PassOptionsToPackage{unicode}{hyperref}
\PassOptionsToPackage{hyphens}{url}
%
\documentclass[
]{article}
\usepackage{amsmath,amssymb}
\usepackage{iftex}
\ifPDFTeX
  \usepackage[T1]{fontenc}
  \usepackage[utf8]{inputenc}
  \usepackage{textcomp} % provide euro and other symbols
\else % if luatex or xetex
  \usepackage{unicode-math} % this also loads fontspec
  \defaultfontfeatures{Scale=MatchLowercase}
  \defaultfontfeatures[\rmfamily]{Ligatures=TeX,Scale=1}
\fi
\usepackage{lmodern}
\ifPDFTeX\else
  % xetex/luatex font selection
\fi
% Use upquote if available, for straight quotes in verbatim environments
\IfFileExists{upquote.sty}{\usepackage{upquote}}{}
\IfFileExists{microtype.sty}{% use microtype if available
  \usepackage[]{microtype}
  \UseMicrotypeSet[protrusion]{basicmath} % disable protrusion for tt fonts
}{}
\makeatletter
\@ifundefined{KOMAClassName}{% if non-KOMA class
  \IfFileExists{parskip.sty}{%
    \usepackage{parskip}
  }{% else
    \setlength{\parindent}{0pt}
    \setlength{\parskip}{6pt plus 2pt minus 1pt}}
}{% if KOMA class
  \KOMAoptions{parskip=half}}
\makeatother
\usepackage{xcolor}
\usepackage[margin=1in]{geometry}
\usepackage{longtable,booktabs,array}
\usepackage{calc} % for calculating minipage widths
% Correct order of tables after \paragraph or \subparagraph
\usepackage{etoolbox}
\makeatletter
\patchcmd\longtable{\par}{\if@noskipsec\mbox{}\fi\par}{}{}
\makeatother
% Allow footnotes in longtable head/foot
\IfFileExists{footnotehyper.sty}{\usepackage{footnotehyper}}{\usepackage{footnote}}
\makesavenoteenv{longtable}
\usepackage{graphicx}
\makeatletter
\newsavebox\pandoc@box
\newcommand*\pandocbounded[1]{% scales image to fit in text height/width
  \sbox\pandoc@box{#1}%
  \Gscale@div\@tempa{\textheight}{\dimexpr\ht\pandoc@box+\dp\pandoc@box\relax}%
  \Gscale@div\@tempb{\linewidth}{\wd\pandoc@box}%
  \ifdim\@tempb\p@<\@tempa\p@\let\@tempa\@tempb\fi% select the smaller of both
  \ifdim\@tempa\p@<\p@\scalebox{\@tempa}{\usebox\pandoc@box}%
  \else\usebox{\pandoc@box}%
  \fi%
}
% Set default figure placement to htbp
\def\fps@figure{htbp}
\makeatother
\setlength{\emergencystretch}{3em} % prevent overfull lines
\providecommand{\tightlist}{%
  \setlength{\itemsep}{0pt}\setlength{\parskip}{0pt}}
\setcounter{secnumdepth}{-\maxdimen} % remove section numbering
\usepackage{bookmark}
\IfFileExists{xurl.sty}{\usepackage{xurl}}{} % add URL line breaks if available
\urlstyle{same}
\hypersetup{
  pdftitle={Final Project Presentation},
  hidelinks,
  pdfcreator={LaTeX via pandoc}}

\title{Final Project Presentation}
\author{}
\date{\vspace{-2.5em}}

\begin{document}
\maketitle

\subsection{Description}\label{description}

Students will select a policy issue of interest and deliver a 10-minute
presentation summarizing the quantitative analysis conducted to address
their research question(s). Presentations should include the following:

\begin{itemize}
\tightlist
\item
  The relevance and significance of the policy issue (include
  appropriate references)
\item
  Research question(s)
\item
  A description of the quantitative analysis:

  \begin{itemize}
  \tightlist
  \item
    What was done and why this approach was selected
  \item
    Why the chosen method is appropriate for answering the research
    question(s)
  \item
    Key findings
  \item
    Interpretation of findings, including their implications for
    policymaking, explained in clear and accessible terms
  \end{itemize}
\end{itemize}

\subsection{Grading rubric:}\label{grading-rubric}

\begin{longtable}[]{@{}
  >{\raggedright\arraybackslash}p{(\linewidth - 4\tabcolsep) * \real{0.2933}}
  >{\raggedright\arraybackslash}p{(\linewidth - 4\tabcolsep) * \real{0.2933}}
  >{\raggedleft\arraybackslash}p{(\linewidth - 4\tabcolsep) * \real{0.3867}}@{}}
\toprule\noalign{}
\begin{minipage}[b]{\linewidth}\raggedright
Criterion
\end{minipage} & \begin{minipage}[b]{\linewidth}\raggedright
Description
\end{minipage} & \begin{minipage}[b]{\linewidth}\raggedleft
Points
\end{minipage} \\
\midrule\noalign{}
\endhead
\bottomrule\noalign{}
\endlastfoot
\textbf{Policy Issue \& Research Question(s)} & Clearly explains the
relevance and significance of the policy issue, supported by appropriate
references, and states clear research question(s). & 4 \\
\textbf{Quantitative Analysis} & Clearly describes what analysis was
conducted, why it was selected, and why the method is appropriate for
answering the research question(s). & 5 \\
\textbf{Key Findings} & Clearly presents and summarizes the key findings
using appropriate quantitative evidence, such as statistics, tables, or
figures. & 4 \\
\textbf{Interpretation \& Policy Implications} & Accurately interprets
the findings and clearly explains their meaning and implications for
policymaking in accessible, non-technical terms. & 4 \\
\textbf{Presentation \& Organization} & Presentation is clear, well
organized, and appropriately paced for the 10-minute time limit. & 3 \\
\textbf{Total} & & \textbf{20} \\
\end{longtable}

\end{document}
